\section{Mapping}\label{sec:mapping}

\paragraph{Gridding} Footprints are placed on a 2\textdegree{} latitude--longitude grid. Each footprint quadrilateral is sampled with a $5\times5$ bilinear grid of points between its four corners, and its weight is shared equally among the grid cells containing those points. The footprint corners recorded in the Level-1 headers span a median of \FootprintB{}\,km across track by \FootprintA{}\,km along track. This is the collimator field of view, about \SI{25}{km} at \SI{100}{km} altitude, plus roughly \SI{13}{km} of ground-track motion during the \SI{8}{s} integration. One degree is \SI{30.3}{km} at the lunar equator, so a footprint overlaps one to four 2\textdegree{} cells.

\paragraph{Cell values and quality layers} Each cell value is the inverse-variance-weighted mean of the footprint ratios. Three quality layers accompany each map:
\begin{itemize}
\item the number of footprints contributing to the cell;
\item the standard error of the weighted mean;
\item the weighted scatter of the footprint ratios.
\end{itemize}

\paragraph{Coverage} We report coverage separately for all dayside footprints and for the quality-0 mapping set. We give it both as a fraction of grid cells and as a fraction of lunar surface area, where cells are weighted by $\cos(\text{latitude})$. No smoothing, interpolation or infill enters any statistic, map value or validation number in this paper.

\paragraph{Resolution} The effective spatial resolution of a single measurement is the footprint, about $25\times37$\,km. The 2\textdegree{} cells (\SI{61}{km} at the equator) set the resolution of the maps. The grid spacing should not be read as a resolution finer than either of these.
